% Part: first-order-logic
% Chapter: introduction
% Section: soundness-completeness

\documentclass[../../../include/open-logic-section]{subfiles}

\begin{document}

\olfileid[hi]{fol}{int}{scp}

\olsection{सुदृढ़ता और पूर्णता}

हम प्रथम-क्रम तर्क के लिए !!{derivation} प्रणालियाँ भी प्रस्तुत करेंगे।
तर्कशास्त्रियों ने अनेक !!{derivation} प्रणालियाँ विकसित की हैं, किंतु वे
सभी !!{sentence}s के बीच वही !!{derivability} संबंध परिभाषित करती हैं। हम
कहते हैं कि $\Gamma$,~$!A$ को \emph{!!{derive}s} करता है, अर्थात्
$\Gamma \Proves !A$, यदि किसी सटीक रूप से परिभाषित प्रकार की
!!a{derivation} हो। !!^{derivation}s सदा चिह्नों की परिमित व्यवस्थाएँ होती
हैं---संभवतः !!{sentence}s की सूची, अथवा कोई अधिक जटिल संरचना।
!!{derivation} प्रणालियों का उद्देश्य यह निर्धारित करने का साधन देना है कि
कोई !!a{sentence}, किसी समुच्चय~$\Gamma$ से निष्पन्न होता है या नहीं। इस
उद्देश्य को पूरा करने के लिए यह सत्य होना चाहिए कि $\Gamma \Entails !A$
तभी और केवल तभी जब $\Gamma
\Proves !A$।

यदि $\Gamma \Proves !A$ किंतु $\Gamma \Entails !A$ नहीं, तो हमारी
!!{derivation} प्रणाली अत्यधिक प्रबल होगी और आवश्यकता से अधिक सिद्ध करेगी।
यदि $\Gamma
\Proves !A$, तो $\Gamma \Entails !A$---इस गुण को
\emph{सुदृढ़ता} कहते हैं, और यह किसी भी अच्छी !!{derivation} प्रणाली की
न्यूनतम अपेक्षा है। दूसरी ओर, यदि $\Gamma \Entails !A$ किंतु
$\Gamma \Proves !A$ नहीं, तो हमारी !!{derivation} प्रणाली अत्यधिक दुर्बल
होगी और पर्याप्त सिद्ध नहीं करेगी। यदि $\Gamma \Entails !A$, तो
$\Gamma \Proves !A$---इस गुण को \emph{पूर्णता} कहते हैं। सुदृढ़ता को सिद्ध
करना सामान्यतः अपेक्षाकृत सरल है (आगमनात्मक रूप से परिभाषित
!!{derivation}s जिस संरचना के अनुसार बनती हैं, उस संरचना पर आगमन द्वारा)।
पूर्णता को सिद्ध करना अधिक कठिन है।

सुदृढ़ता और पूर्णता के अनेक महत्त्वपूर्ण परिणाम हैं। यदि !!{sentence}s का
कोई समुच्चय~$\Gamma$ किसी विरोधाभास (जैसे~$!A \land \lnot !A$) को
!!{derive}s करता है, तो उसे \emph{असंगत} कहते हैं। असंगत~$\Gamma$ का कोई
प्रतिरूप नहीं हो सकता; वह असंतुष्टनीय होता है। पूर्णता से इसका विलोम मिलता है:
जो भी~$\Gamma$ असंगत नहीं है---अर्थात्, जैसा हम कहेंगे, \emph{संगत} है---उसका
एक प्रतिरूप होता है। वास्तव में यह पूर्णता के समतुल्य है और पूर्णता का वही
रूप है जिसे हम सचमुच सिद्ध करेंगे। यह एक गहन और संभवतः आश्चर्यजनक परिणाम है:
केवल यह कि आप $!A \land \lnot !A$ को~$\Gamma$ से सिद्ध नहीं कर सकते, इस बात
की गारंटी देता है कि ऐसी !!a{structure} है जो~$\Gamma$ द्वारा किए गए वर्णन
के अनुरूप है। अतः पूर्णता इस प्रश्न का उत्तर देती है: !!{sentence}s के किन
समुच्चयों के प्रतिरूप होते हैं? उत्तर: ठीक सभी संगत समुच्चयों के।

सुदृढ़ता और पूर्णता के दो महत्त्वपूर्ण परिणाम हैं। पहला संहतता प्रमेय है और
दूसरा लोवेनहाइम--स्कोलेम (L\"owenheim--Skolem) प्रमेय। ये प्रतिरूप सिद्धान्त के
महत्त्वपूर्ण परिणाम हैं और अनेक रोचक परिणाम स्थापित करने में उपयोग किए जा
सकते हैं। हम पहले ही दो का उल्लेख कर चुके हैं: प्रथम-क्रम तर्क यह व्यक्त
नहीं कर सकता कि !!a{structure} का !!{domain} परिमित है अथवा वह
!!{nonenumerable} है।

ऐतिहासिक रूप से यह सब समझने और ठीक करने में बहुत समय लगा---प्रथम-क्रम तर्क
का वाक्यविन्यास और अर्थविज्ञान कैसे परिभाषित करें, अच्छी !!{derivation}
प्रणालियाँ कैसे परिभाषित करें, उन्हें सुदृढ़ और पूर्ण कैसे सिद्ध करें, और
प्रथम-क्रम भाषाओं में क्या व्यक्त किया जा सकता है तथा क्या नहीं, इसे स्पष्ट
कैसे करें। अब हम जानते हैं कि यह कैसे करना है, किंतु सभी विवरणों से गुजरना
फिर भी उलझाने वाला और श्रमसाध्य हो सकता है। फिर भी यह महत्त्वपूर्ण है, क्योंकि
प्रथम-क्रम तर्क की औपचारिक भाषा के लिए यहाँ विकसित विधियाँ तर्कशास्त्र,
संगणक विज्ञान और भाषाविज्ञान में व्यापक रूप से लागू होती हैं। अतः विवरणों पर
मेहनत करना अंततः फलदायी होता है।

\end{document}
